\begin{problem}[第32届全国中学生物理竞赛复赛][核聚变循环]{40}
    在太阳内部存在两个主要的核聚变反应过程：碳循环和质子-质子循环；其中碳循环是贝蒂在1938年提出的，碳循环反应过程如图所示。图中 $\mathrm{p}$、$\mathrm{e}^{+}$ 和 $\nu_{\mathrm{e}}$ 分别表示质子、正电子和电子型中微子；粗箭头表示循环反应进行的先后次序。当从循环图顶端开始，质子 $\mathrm{p}$ 与 $^{12}\mathrm{C}$ 核发生反应生成 $^{13}\mathrm{N}$ 核，反应按粗箭头所示的次序进行，直到完成一个循环后，重新开始下一个循环。已知 $\mathrm{e}^{+}$、$\mathrm{p}$ 和 $\mathrm{He}$ 核的质量分别为 $0.511\,\mathrm{MeV}/c^{2}$、$1.0078\,\mathrm{u}$ 和 $4.0026\,\mathrm{u}$（$1\,\mathrm{u}\approx 931.494\,\mathrm{MeV}/c^{2}$），电子型中微子 $\nu_{\mathrm{e}}$ 的质量可以忽略。
    \begin{subproblem}
        写出图中 X 和 Y 代表的核素；
    \end{subproblem}
    \begin{subproblem}
        写出一个碳循环所有的核反应方程式；
    \end{subproblem}
    \begin{subproblem}
        计算完成一个碳循环过程释放的核能。
    \end{subproblem}
\end{problem}